\chapter{مقدمه‌ای بر داده‌کاوی متن و وب}
\label{ch:text-web}

\section*{اهداف این فصل}
\begin{itemize}
    \item آشنایی با گام‌های اصلی پیش‌پردازش متن
    \item یادگیری بازنمایی متن به‌صورت عددی (Bag-of-Words و TF-IDF)
    \item آشنایی مقدماتی با کاربرد طبقه‌بندی و خوشه‌بندی روی متن
    \item آشنایی کلی با مفاهیم پایهٔ داده‌کاوی وب
\end{itemize}

\section{چرا متن به رفتار متفاوتی نیاز دارد؟}

تمام الگوریتم‌هایی که تاکنون دیدیم (درخت تصمیم، k-means، k-NN و
مانند آن‌ها) به یک بازنمایی عددی و با ابعاد ثابت از هر نمونه نیاز
دارند. اما متن خام — یک ایمیل، یک نظر مشتری، یک مقالهٔ خبری — رشته‌ای
از کلمات با طول متغیر است که مستقیماً قابل استفاده در این الگوریتم‌ها
نیست. بنابراین اولین گام در هر پروژهٔ
\dmterm{\textbf{داده‌کاوی متن}}{Text Mining}، تبدیل متن به یک بازنمایی
عددی مناسب است.

\section{پیش‌پردازش متن}

پیش از تبدیل متن به بردار عددی، معمولاً چند گام پیش‌پردازش زیر انجام
می‌شود:

\begin{itemize}
    \item \dmterm{\textbf{نشانه‌گذاری}}{Tokenization}: شکستن متن به
          واحدهای کوچک‌تر، معمولاً کلمات یا زیرکلمات\footnote{\LR{Tokens}}.
    \item \dmterm{\textbf{حذف ایست‌واژه‌ها}}{Stop-word Removal}: حذف
          کلمات بسیار پرتکرار و کم‌اطلاعات مانند حروف اضافه و ضمایر
          (مثلاً «را»، «به»، «از» در فارسی، یا «the»، «is» در انگلیسی)
          که معمولاً برای تشخیص موضوع متن مفید نیستند.
    \item \dmterm{\textbf{ریشه‌یابی}}{Stemming} یا
          \dmterm{\textbf{بن‌واژه‌یابی}}{Lemmatization}: تبدیل کلمات
          مختلف با ریشهٔ مشترک به یک شکل واحد (مثلاً «رفتن»، «می‌رود»،
          «رفت» به یک ریشهٔ مشترک)، تا الگوریتم آن‌ها را یک مفهوم واحد
          در نظر بگیرد.
    \item نرمال‌سازی متن: تبدیل به حروف کوچک (در زبان‌های لاتین)، حذف
          علائم نگارشی و اعداد نامرتبط.
\end{itemize}

\section{بازنمایی متن به‌صورت عددی}

\subsection{کیسهٔ کلمات}

ساده‌ترین روش بازنمایی متن،
\dmterm{\textbf{کیسهٔ کلمات}}{Bag-of-Words (BoW)} است: هر سند با یک
بردار به طول واژگان کل مجموعه داده نمایش داده می‌شود که هر درایهٔ آن،
تعداد تکرار یک کلمهٔ خاص در آن سند است. نام «کیسه» به این دلیل است که
این روش کاملاً \textbf{ترتیب} کلمات را نادیده می‌گیرد؛ برای این روش،
جملهٔ «سگ گربه را دنبال کرد» و «گربه سگ را دنبال کرد» کاملاً یکسان
دیده می‌شوند.

\subsection{TF-IDF}

مشکل اصلی کیسهٔ کلمات این است که به کلمات بسیار پرتکرار (که لزوماً
اطلاعات زیادی دربارهٔ موضوع سند نمی‌دهند) وزن بیش از حد می‌دهد.
\dmterm{\textbf{\LR{TF-IDF}}}{Term Frequency-Inverse Document Frequency}
این مشکل را با ترکیب دو عامل حل می‌کند:

\begin{itemize}
    \item \dmterm{\textbf{فراوانی واژه}}{Term Frequency (TF)}: تعداد
          تکرار یک کلمه در یک سند خاص (هرچه بیشتر، وزن بیشتر).
    \item \dmterm{\textbf{فراوانی معکوس سند}}{Inverse Document Frequency
          (IDF)}: هرچه یک کلمه در تعداد \textbf{کمتری} از اسناد کل
          مجموعه ظاهر شود، وزن بیشتری می‌گیرد؛ کلماتی که تقریباً در همهٔ
          اسناد تکرار می‌شوند (مانند حروف اضافهٔ باقی‌مانده)، وزن بسیار
          کمی می‌گیرند.
\end{itemize}

وزن نهایی هر کلمه در هر سند، حاصل‌ضرب این دو مقدار است:
\[
\text{TF-IDF}(t, d) = \text{TF}(t, d) \times \text{IDF}(t)
\]
این بازنمایی، پرکاربردترین روش کلاسیک برای تبدیل متن به بردار عددی است
و پایهٔ بسیاری از موتورهای جستجو و سیستم‌های بازیابی اطلاعات
نیز همین ایده است\footnote{\LR{Information Retrieval}}.

\section{کاربردهای مقدماتی}

پس از بازنمایی متن به‌صورت بردار عددی (با \LR{TF-IDF} یا کیسهٔ کلمات)،
می‌توان مستقیماً از الگوریتم‌های فصل‌های قبل استفاده کرد:

\begin{itemize}
    \item \textbf{طبقه‌بندی متن}: مثلاً تشخیص اسپم یا تحلیل احساسات
          (مثبت/منفی/خنثی بودن یک نظر مشتری) با استفاده از بیز ساده یا
          ماشین بردار پشتیبان.
    \item \textbf{خوشه‌بندی متن}: گروه‌بندی خودکار مقالات خبری بر اساس
          موضوع، بدون داشتن برچسب از پیش، با استفاده از k-means روی
          بردارهای \LR{TF-IDF}.
\end{itemize}

نکتهٔ مهم این است که بردارهای حاصل از متن معمولاً \textbf{ابعاد بسیار
بالا} (به تعداد کل کلمات منحصربه‌فرد واژگان) و \textbf{بسیار پراکنده}
دارند (اکثر درایه‌ها صفرند، چون هر سند فقط شامل بخش کوچکی از کل واژگان
است)؛ به همین دلیل، همان‌طور که در فصل \ref{ch:data} دیدیم، شباهت
کسینوسی معمولاً معیار مناسب‌تری نسبت به فاصلهٔ اقلیدسی برای مقایسهٔ این
بردارهاست.

\section{مقدمه‌ای بر داده‌کاوی وب}

\dmterm{\textbf{داده‌کاوی وب}}{Web Mining} به کاربرد تکنیک‌های
داده‌کاوی بر روی داده‌های وب اشاره دارد و معمولاً به سه زیرشاخه تقسیم
می‌شود:

\begin{itemize}
    \item \dmterm{\textbf{کاوش محتوای وب}}{Web Content Mining}: استخراج
          دانش از محتوای صفحات وب (متن، تصویر)، که عملاً تعمیمی از
          داده‌کاوی متن به مقیاس وب است.
    \item \dmterm{\textbf{کاوش ساختار وب}}{Web Structure Mining}: تحلیل
          گراف پیوندهای میان صفحات وب (مثلاً الگوریتم‌هایی شبیه
          \LR{PageRank} که برای رتبه‌بندی صفحات از ساختار پیوندها
          استفاده می‌کنند).
    \item \dmterm{\textbf{کاوش استفاده از وب}}{Web Usage Mining}: تحلیل
          رفتار کاربران در یک وب‌سایت (مانند مسیر کلیک‌ها) برای بهبود
          طراحی سایت یا سیستم‌های توصیه‌گر.
\end{itemize}

پرداختن کامل به این حوزه از حوصلهٔ این کتاب مقدماتی خارج است، اما آشنایی
با این تقسیم‌بندی برای شناخت جایگاه داده‌کاوی متن در تصویر بزرگ‌تر
مفید است.

\section*{پیاده‌سازی در پایتون}

پیاده‌سازی \texttt{TfidfVectorizer} و \texttt{CountVectorizer} از
\texttt{scikit-learn}، به‌همراه یک مثال کامل طبقه‌بندی متن، در بخش
«فصل ۹» پیوست الف (\ref{app:python}) آمده است.

\section*{خلاصهٔ فصل}

در این فصل با چالش اصلی کار با داده‌های متنی — نیاز به تبدیل متن به
بازنمایی عددی — آشنا شدیم. گام‌های پیش‌پردازش متن، کیسهٔ کلمات و
\LR{TF-IDF} را بررسی کردیم و دیدیم چگونه پس از این تبدیل، می‌توان از
همان الگوریتم‌های طبقه‌بندی و خوشه‌بندی فصل‌های قبل روی متن استفاده
کرد. در پایان، نگاهی کلی به سه زیرشاخهٔ داده‌کاوی وب انداختیم. در فصل
پایانی کتاب، به‌طور مقدماتی با داده‌های پیچیده‌تر (گراف و سری زمانی)
آشنا خواهیم شد.

\section*{تمرین‌ها}

\begin{enumerate}
    \item چرا نمی‌توان متن خام را مستقیماً به یک الگوریتم مانند درخت
          تصمیم یا k-means داد؟
    \item دو نمونه از ایست‌واژه‌های فارسی را نام ببرید و توضیح دهید چرا
          حذف آن‌ها معمولاً به بهبود کیفیت مدل کمک می‌کند.
    \item فرض کنید کلمهٔ «کامپیوتر» در یک سند خاص ۵ بار تکرار شده، اما
          در تقریباً همهٔ اسناد مجموعه نیز حضور دارد. وزن \LR{TF-IDF}
          این کلمه چه تفاوتی با کلمه‌ای کمیاب‌تر اما با همان فراوانی
          خواهد داشت؟
    \item چرا شباهت کسینوسی برای مقایسهٔ بردارهای \LR{TF-IDF} مناسب‌تر
          از فاصلهٔ اقلیدسی است؟
    \item تفاوت کاوش محتوای وب و کاوش ساختار وب را با یک مثال از هرکدام
          توضیح دهید.
\end{enumerate}
